% TestVDB v2 — target outline (Stage 1, awaiting approval).
% Spine = novelty-reassessment.md §8/§9. Body content reused from paper-draft-vldb-final.tex.
% Format: acmart sigconf (double-column ACM SIG). Page limit: none (venue TBD).
% DO NOT write section prose until the user approves this outline.
\documentclass[sigconf]{acmart}

\settopmatter{printacmref=false}   % venue TBD; suppress boilerplate for now
\settopmatter{printfolios=false}

\title[TestVDB: source-grounded falsification for VDB API compliance]{TestVDB: Source-Grounded Falsification of LLM-Derived Contracts for API-Compliance Testing of Vector Databases}

\author{Authors TBD}
\email{tbd@example.org}
\affiliation{\institution{TBD (anonymity depends on venue)}\country{TBD}}

\begin{abstract}
% Outline (rewrite per §8/§9, NOT the old propagation-centric abstract):
% 1. Problem: VDB API compliance defects (~43\% of VDB bugs are incorrect-behavior) lack practical oracles; fuzzers catch only crash.
% 2. Finding (E1): ~90\% of real VDB compliance defects form a residual that classical oracles (differential / metamorphic / property-based) and REST doc$\to$assertion oracles (AGORA+/SATORI/MASTOR) structurally cannot reach, because accept/reject semantics are not compilable to a deterministic assertion --- forcing an LLM checker (the \emph{LLM-as-checker regime}).
% 3. Finding (t25): in that regime, contract errors are partly \emph{task-intrinsic} (all LLM families reproduce them from ambiguous docs); cross-model validation fixes only the family-specific subset (Panickssery-style self-preference).
% 4. Method: TestVDB --- source-grounded falsification that resolves BOTH subsets using implementation as ground truth; complementary to cross-model, not redundant.
% 5. Results: 5 VDBs, 111 defects submitted, 38 acknowledged; source anchor lifts FP-suppression 31\%$\to$81\% at 96.7\% TP retention.
[Abstract to be written after outline approval.]
\end{abstract}

\begin{CCSXML}
<ccs2012><concept><concept_id>10011007.10011074.10011099</concept_id><concept_desc>Software and its engineering~Software testing and debugging</concept_desc><concept_significance>500</concept_significance></concept></ccs2012>
\end{CCSXML}
\ccsdesc[500]{Software and its engineering~Software testing and debugging}

\keywords{Vector database, API compliance testing, test oracle, LLM-as-judge, source-grounded falsification, conformance testing}

\begin{document}

\maketitle

% =====================================================================
\section{Introduction}
\textit{Plan (per §8/§9).}
\begin{itemize}
  \item VDBMS underpin LLM apps; $\sim$43\% of VDB bugs are incorrect-behavior, not crash; fuzzers (VDBFuzz) cover only crash. We target \emph{API compliance defects}: silent accept/behavior violating the documented contract.
  \item \textbf{Empirical anchor (E1):} across our 111 defects, $\sim$90\% are a compliance residual that differential/metamorphic/PBT cannot reach (they cover the math+crash subset); REST doc$\to$assertion oracles (AGORA+/SATORI/MASTOR) keep a deterministic checker and so never enter the regime where compliance lives.
  \item \textbf{The LLM-as-checker regime:} compliance accept/reject is not compilable to a deterministic assertion $\Rightarrow$ the checker is forced to be an LLM. This is the structural reason an LLM is needed --- not a fashion choice.
  \item \textbf{The reliability problem inside the regime:} self-preference (Panickssery et al. 2024, family-specific) AND task-intrinsic contract errors (t25: cross-model fixes only the family-specific half; the task-intrinsic half needs implementation ground truth).
  \item \textbf{Our answer:} TestVDB --- source-grounded falsification that resolves both subsets; complementary to cross-model. CTS (Contract-Truth Separation) is the mechanism.
  \item \textbf{Contributions:} (C1) the LLM-as-checker-regime framing + E1 quantification of the compliance residual; (C2) source-grounded falsification as the task-intrinsic resolver (t25 evidence); (C3) system + 111 defects / 38 acknowledged across 5 VDBs; (C4) a reusable model-free invariant oracle subclass.
  \item Exclusion Table 1 (already redone in v1) goes here.
\end{itemize}

% =====================================================================
\section{Background and Problem Setup}
\textit{Plan.}
\begin{itemize}
  \item Oracle taxonomy (Barr 2015): specified / derivable / implicit / none $\Rightarrow$ we sit in specified (contract) + derivable (math invariant), and argue the compliance residual is where all four classical classes fail.
  \item VDB compliance vs correctness: compliance = accept/reject vs documented contract; correctness = result math (ANN recall). We do compliance; correctness stays open.
  \item Why compliance is non-compilable: the accepted set is defined by NL docs, no formal grammar $\Rightarrow$ cannot become a deterministic assertion. (Sets up §3.)
  \item Reuse bug-distribution framing (roadmap25 / bugstudy25) from v1.
\end{itemize}

% =====================================================================
\section{The LLM-as-Checker Regime} \label{sec:regime}
\textit{Plan --- NEW section, the conceptual core (§8/§9). This is what v1 lacked.}
\begin{itemize}
  \item \textbf{Definition:} a setting is in the LLM-as-checker regime when the oracle cannot be compiled to a deterministic assertion, so the pass/fail judgment must come from an LLM.
  \item \textbf{Why VDB compliance forces it:} accept/reject semantics are NL-defined; contrast with AGORA+/SATORI/MASTOR, whose oracles ARE executable assertions $\Rightarrow$ deterministic checker $\Rightarrow$ they stay OUT of this regime (this is the precise distinction from the REST oracle line).
  \item \textbf{Two failure modes inside the regime:}
    \begin{itemize}
      \item \emph{Family-specific} (self-preference, Panickssery 2024): same family generates contract + judges $\Rightarrow$ shared bias self-confirms. Mitigated by cross-model.
      \item \emph{Task-intrinsic}: ambiguous docs make ALL families infer the same wrong (e.g., over-strict) contract. Cross-model CANNOT fix this --- both agree wrongly.
    \end{itemize}
  \item \textbf{t25 evidence:} feed the same doc passage to GLM vs DeepSeek; DeepSeek reproduces GLM's over-strict contract in only $\sim$half of cases $\Rightarrow$ the other half is task-intrinsic. (t25 N=9 pilot; expand for camera-ready.)
  \item \textbf{Implication:} source-grounded falsification is needed precisely for the task-intrinsic subset --- implementation is the only ground truth that resolves what cross-model cannot. This is the precise, defensible ``why source'' argument.
\end{itemize}

% =====================================================================
\section{TestVDB Design} \label{sec:design}
\textit{Plan --- reuse v1 mechanism, reframe CTS as the task-intrinsic resolver.}
\begin{itemize}
  \item 5-stage pipeline, 20 agents (reuse v1 §3): contract extraction $\to$ attack generation $\to$ 4-judge debate $\to$ dev-reviewer (source-grounded falsification) $\to$ novelty gate.
  \item Contract-Truth Separation (CTS): assertion layer (LLM) vs truth layer (source-grounded falsifier) --- reframe as ``the truth layer resolves task-intrinsic contract errors that cross-model cannot.''
  \item dev-reviewer's three anchors: clean reproduction, source-grounded verification (primary), threat-model cross-check (optional).
  \item Model-free invariant oracle subclass (COSINE$\in[-1,1]$, index completeness) --- the classical-addressable slice, reusable standalone.
\end{itemize}

% =====================================================================
\section{Implementation} \label{sec:impl}
\textit{Plan --- reuse v1. GLM-5.2 backbone via Claude Code runtime; 20 agent prompts; target version matrix; cost ($\sim$10/target, $10^4$ calls).}

% =====================================================================
\section{Evaluation} \label{sec:eval}
\textit{Plan --- reuse v1 eval structure, elevate E1 and t25 to first-class RQs.}
\begin{itemize}
  \item \textbf{RQ1 (yield + E1):} 111 submissions, 38 acknowledged across 5 VDBs; \emph{fault-model split} (E1): $\sim$90\% compliance residual, $\sim$10\% classical-addressable (math+crash). Table of per-DB yield + the M/C/X/V breakdown.
  \item \textbf{RQ2 (source falsification effect):} retrospective 52 candidates; source anchor lifts FP-suppression 31\%$\to$81\%, TP retention 96.7\%; e2e precision 69.2\% (Wilson CI + pending bound).
  \item \textbf{RQ3 (regime evidence / t25):} cross-family counterfactual --- GLM vs DeepSeek contract formalization on the same docs; family-specific vs task-intrinsic split. (Currently t25 N=9; state as pilot + camera-ready expansion plan.)
  \item \textbf{RQ4 (model-free invariant):} COSINE$>$1 etc., cross-vendor (Milvus+Qdrant), the classical-addressable subclass.
  \item Threats to validity: t25 N=9; single model family for CTS; soft result-correctness open; breadth-only for 3 VDBs.
\end{itemize}

% =====================================================================
\section{Related Work} \label{sec:related}
\textit{Plan --- restructure around the regime distinction (closes the v1 zero-citation hole).}
\begin{itemize}
  \item \textbf{VDBMS testing:} VDBFuzz (crash), roadmap25, bugstudy25.
  \item \textbf{REST doc$\to$oracle line:} AGORA+ (Daikon dynamic invariants), SATORI (LLM from OpenAPI), MASTOR (multi-agent from source). Position via the deterministic-checker regime: all three stay out of the LLM-as-checker regime; MASTOR additionally treats code as truth (mutation paradigm) and explicitly cannot detect doc-vs-code compliance (§7.4.2) --- orthogonal fault model.
  \item \textbf{LLM-as-judge reliability:} Panickssery et al. 2024 (self-preference bias) --- anchor for the family-specific subset; we extend to the pipeline setting + identify the task-intrinsic residual beyond it.
  \item \textbf{DB correctness oracles:} NoREC/TLP/DQE/DDLCheck (SQL semantics; not applicable to VDB compliance).
  \item \textbf{Property-based / metamorphic:} QuickCheck, Schemathesis, QuickREST, MeTMaP --- cover the math subset, structurally miss compliance (Table 1).
\end{itemize}

% =====================================================================
\section{Discussion and Limitations}
\textit{Plan.}
\begin{itemize}
  \item t25 is a pilot (N=9); camera-ready expansion via live special-value testing across 5 VDBs.
  \item Cross-model CTS validation (full 3-arm) is the primary open direction.
  \item Soft result-correctness (ANN recall/ranking) remains open.
  \item Breadth-only generalization for Weaviate/MeiliSearch/Chroma; statistical claims on Milvus+Qdrant.
  \item Terminology: we adopt \emph{conformance} (NIST) over \emph{compliance} where precision matters (Task-3 audit).
\end{itemize}

% =====================================================================
\section{Conclusion}
\textit{Plan --- one paragraph: VDB API compliance is a residual that forces the LLM-as-checker regime; the regime's task-intrinsic contract errors are resolvable only by source-grounded falsification (complementary to cross-model); TestVDB demonstrates at scale (111 defects, 38 acknowledged).}

% =====================================================================
% Outline ends. Awaiting user approval before writing section prose.
\end{document}
